\centerline{\bf \Large  РЕШЕНИЯ и КРИТЕРИИ}

\centerline{\bf \Large  8 КЛАСС}


\problem{\textbf{1}} Выпишем все возможные варианты трёх натуральных чисел, произведение которых равно 36.

$36\cdot1\cdot1=36$\\
$18\cdot2\cdot1=36$\\
$12\cdot3\cdot1=36$\\
$9\cdot4\cdot1=36$\\
$9\cdot2\cdot2=36$\\
$6\cdot6\cdot1=36$\\
$6\cdot3\cdot2=36$\\
$4\cdot3\cdot3=36$\\

Два числа в произведении можем делать отрицательными и перебором доказываем, что среди этих вариантов не будет такого, что сумма будет равняться нулю.

\par

\textbf{Критерии.} \\
\textit{7 баллов} $-$ Полностью верное решение через полный перебор\\
\textit{4-6 баллов} $-$ Верное решение с неполным перебором\\
\textit{0-3 балла} $-$  За все остальные попытки решить задачу\\


\problem{\textbf{2}} Пусть за сутки $Y$ клубники съедают насекомые, а $X$ клубники собирает один человек.

Тогда имеем равенство:

$5\cdot(401\cdot X+Y) = 7\cdot(285\cdot X+Y)$

Из которого получаем: $Y=5X$

Всего клубники растёт в огороде $2030X$

Значит, чтобы собрать за один день нужно позвать 2024 друга

$1\cdot((2024+1)\cdot X+5\cdot X) = 2030X$
\par

\textbf{Критерии.} \\
\textit{0-4 баллов} $-$ Решение пункта А\\
\textit{0-1 балла} $-$ Составление условия пункта Б \\
\textit{0-2 баллов} $-$ Решение пункта Б\\
\textit{0 баллов} $-$ Отсутствие решения пункта А (сразу НОЛЬ! баллов)\\
\textit{0 баллов} $-$ Любой ответ без объяснений\\

\problem{\textbf{3}}
Пусть высота $CH$ пересекает отрезок $PQ$ в точке $M$. Обозначим угол $A$ за $\alpha$. Тогда $\angle ABC = \angle ACH =  90 - \alpha , \ \angle A = BCH = \alpha $. Так как четырёхугольник $ABQP$ вписанный, имеем что $\angle CAB = \angle PQC = \alpha$  и $\angle CBA = \angle QPC = 90 - \alpha$. Стало быть треугольники $PMC$ и $CMQ$ равнобедренные причём $PM = MC = MQ$. То есть $M$ середина $PQ$.

\textbf{Критерии.} \\
\textit{7 баллов} $-$ Полностью верно решение подсчётом углов\\
\textit{7 баллов} $-$ Решение использующее свойства симедианы\\
\textit{1 балл} $-$  Показано, что $\angle ABC = \angle ACH , \ \angle CAB = BCH$ \\ 
\textit{1 балл} $-$ Равенство углов при вписанном четырёхугольнике\\


\problem{\textbf{4}} Оценка: Отметим все числа с 1000000 до 10000000 с шагом в 1000. Таким образом, мы отметили все НЕУДЕкрасивые числа. Значит, УДЕкрасивых чисел может быть не большее 999, так как шаг в 1000.

Пример: Самое маленькое НЕУДЕкрасивое число - это $1000\cdot1000=1000000$, следующее же после него НЕУДЕкрасивое число - это $1000\cdot1001=1001000$. Между ними будут 999 УДЕкрасивых чисел.

1000001, 1000002, ... 1000999\par

\textbf{Критерии.} \\
\textit{7 баллов} $-$ Полностью верное решение\\
\textit{5 баллов} $-$ Верная оценка\\
\textit{2 балла} $-$ Верный пример с объяснением через минимальные НЕУДЕкрасивые числа\\
\textit{1 балл} $-$  Верный пример чисел с 1000001, 1000002, ... 1000999 без объяснения того, что между ними не может быть УДЕкрасивых\\ 
\textit{0-1 балл} $-$ За все остальные попытки решить задачу\\


\problem{\textbf{5}}
Если борцоки брал Арслан, в коробке остается 16/17, если Данзан - 11/12, если Очир - $3 / 4$ от их предыдущего количества. Если Арслан, Дазнан и Очир брали борцоки из данной коробки несколько раз, их первоначальное количество умножается на 16/17, 11/12 и 3/4 соответствующее количество раз. Если сначала борцоков в обеих коробках было поровну, и в конце - тоже, то дроби, на которые умножалось количество борцоков в первой и второй коробках, также равны. Но тогда у них в несократимой записи в числителях должно быть поровну множителей 11 и поровну множителей 3 , а в знаменателях - поровну множителей 17. А так как всем этим множителям не с чем сокращаться, то их было поровну и в исходной записи, откуда и следует справедливость утверждения задачи.

\textbf{Критерии.} \\
\textit{7 баллов} $-$ Полностью верное решение \\
\textit{3 балла} $-$ Упомянута основная теорема арифметики, как инструмент решения задачи\\
\textit{2 балла} $-$ Показано, что при взятии борцоков тем или иным мальчиком количество борцоков умножается на некоторую дробь\\
\textit{1 балл} $-$  Доказано, утверждение только для двух мальчиков \\ 
\textit{0 баллов} $-$ За все остальные попытки решить задачу\\